\subsection{前端设计}

前端采用取基本块结构，主要包括分支预测器（BPU）、取指目标队列（FTQ）、取指单元（IFU，含预译码）与一级指令缓存（ICache）。BPU 每拍最多对一个起始地址产生一个预测块，块内最多 4 条指令；出于地址翻译正确性考虑，保证块内指令不跨 4\,KB 页。

\subsubsection{分支预测器}

BPU 实现条件分支、函数调用与无条件跳转的目标预测、取指宽度截断以及条件分支方向预测。核心组件包括：

\begin{itemize}
  \item \textbf{uBTB}：16 项全相联，当拍返回，提供无空泡基线预测块（默认取指宽度 4）；
  \item \textbf{fallback BTB}：64 项直接映射，在 uBTB 未命中时提供 P0 后备目标，避免 CAM （ Content-Addressable Memory ）拉长关键路径；
  \item \textbf{FTB}：四路组相连、每路 2048 组，共 8192 项，使用 BRAM 实现，保存基本块分支元数据；
  \item \textbf{TAGE}：基础表 8192 项（2 bit 饱和计数）与 4 张标记表（各 1024 项，12 bit tag），全局历史长度分别为 11/23/53/112；
  \item \textbf{RAS}：32 项，前端推测栈与后端提交栈双栈结构，分别服务预测期压弹与提交期校正。
\end{itemize}

\textbf{BPU 流水（P0/P1）}：

\begin{itemize}
  \item \textbf{P0（当拍）}：以当前 PC 查询 uBTB / fallback BTB，生成默认预测块并送入 FTQ，保证取指不因主预测延迟而空泡；
  \item \textbf{P1（次拍）}：FTB 与 TAGE 返回。若主预测与 P0 块不一致，则发出前端重定向并修正 FTQ 中上一拍条目。
\end{itemize}

\textbf{FTB 条目}主要包括：由预测块起始地址决定的 tag、顺序下落地址（fall-through）、跳转目标、分支类型等。
FTB/TAGE 的训练以提交级分支结果为准，经训练队列回写，保证推测路径不会污染架构正确的预测表。
当基本块出现无条件跳转、条件分支首次 taken、或分支信息错误等情况时，更新或分配 FTB 条目。

\subsubsection{取指目标队列 FTQ}

FTQ 深度为 16，采用三指针环形结构（BPU 写指针、IFU 读指针、提交释放指针），在分支预测与取指之间解耦，并缓存 BPU 训练元数据。
BPU 在 P0 即将预测块写入 FTQ；若 P1 主预测纠正，则就地修改对应条目。
提交阶段按程序序释放 FTQ 项，并把分支结果送回 BPU 训练。每个块最多 4 条指令，与后端 ROB 深度相匹配，避免前端过度推测导致资源浪费。

\subsubsection{取指令单元与取指流水}

IFU 从 FTQ 读取预测块，经 MMU 做取指地址翻译；无异常时向 ICache 发起整行取指请求。取指流水可概括为：

\begin{itemize}
  \item \textbf{PRE}：接收 FTQ 块，发起虚地址翻译，并在允许时向 ICache 发请求；支持 FTQ$\rightarrow$Icache 同拍直发（在无需主 TLB 的直接翻译场景下）；
  \item \textbf{IF}：等待 ICache 返回 cache line，按预测块起止切割指令，写入 IB；
  \item \textbf{预译码重定向}：对不依赖寄存器的强制跳转做前端检查，若发现漏预测则冲刷前端取指/预测流水并重定向；
  \item \textbf{冲刷协同}：后端冲刷时，根据是否存在在途 Icache 读事务决定是否丢弃迟到返回，并清空流水寄存器。
\end{itemize}

ICache 失效（如 \texttt{cacop}）时，IFU 同步清理 PRE/IF 与行缓冲，避免自修改代码路径读到陈旧指令。
